\documentclass[11pt]{article}
\usepackage[margin=1in]{geometry}
\usepackage{amsmath,amssymb}
\usepackage{enumitem}
\usepackage{bookmark}
\hypersetup{
  pdftitle={PSET 1: Notation, sets and functions, quadratics, and logarithms},
  hidelinks}

\title{PSET 1: Notation, sets and functions, quadratics, and logarithms}
\author{}
\date{\vspace{-2.5em}}

% Problems are numbered 1-7 across the document.
\newcounter{problem}
\newcommand{\problem}[1]{\stepcounter{problem}\section*{Problem \theproblem: #1}}

\begin{document}
\maketitle

% Shared assignment questions (header block).
% Include at the top of any problem set with:  % Shared assignment questions (header block).
% Include at the top of any problem set with:  % Shared assignment questions (header block).
% Include at the top of any problem set with:  \input{assignment-qs}
\section*{Assignment Qs}

\begin{itemize}
\item Name
\item How long did this problem set take you (round to nearest hour)? \quad
  0-1 hour \ 2-3 hours \ 4-5 hours \ 5+ hours
\item How difficult was this problem set? \quad
  very easy \ 1 \ 2 \ 3 \ 4 \ 5 \ very challenging
\end{itemize}

\section*{Assignment Qs}

\begin{itemize}
\item Name
\item How long did this problem set take you (round to nearest hour)? \quad
  0-1 hour \ 2-3 hours \ 4-5 hours \ 5+ hours
\item How difficult was this problem set? \quad
  very easy \ 1 \ 2 \ 3 \ 4 \ 5 \ very challenging
\end{itemize}

\section*{Assignment Qs}

\begin{itemize}
\item Name
\item How long did this problem set take you (round to nearest hour)? \quad
  0-1 hour \ 2-3 hours \ 4-5 hours \ 5+ hours
\item How difficult was this problem set? \quad
  very easy \ 1 \ 2 \ 3 \ 4 \ 5 \ very challenging
\end{itemize}


\problem{Translate a statement into notation}

A researcher proposes a simple model of protest participation: an
individual joins a protest when the benefit they expect from the
protest succeeding, weighted by the chance their participation
matters, exceeds the personal cost of participating (time, risk of
arrest, etc.).

\begin{enumerate}[label=\alph*.]
\item Define notation for each quantity in the model (in the style of
  \(R = PB - C\) from the rational voter model in class) and write the
  condition under which an individual participates.
\item Using your notation, write the formal statement of the following
  implication: ``in a very large protest movement, any one person's
  participation is unlikely to matter, so individuals should stay
  home.''
\item Suppose the government announces harsher penalties for
  protesters. Which term in your model changes, in which direction,
  and what does the model predict about participation?
\end{enumerate}

\problem{Work with sets}

Using the sets\footnote{inspired by Grimmer HW1.1}
\[
\begin{aligned}
A &= \left\{ 2, 3, 7, 9, 12, 13 \right\} \\
B &= \left\{ x : 6 \leq x \leq 10 \ \mbox{and} \ x \ \mbox{is an even integer} \right\} \\
C &= \left\{ x : 2 < x < 20 \ \mbox{and} \ x \ \mbox{is prime} \right\} \\
D &= \left\{ 1, 4, 9, 16, 25, \ldots \right\} \\
E &= \left\{ 1, 4 \right\}
\end{aligned}
\]
identify the following:

\begin{enumerate}[label=\alph*.]
\item \(A \cup B\)
\item \((A \cup B) \cap C\)
\item \(C \cap D\)
\item \(A \times E\)
\end{enumerate}

\problem{Functions vs.\ relations}

\begin{enumerate}[label=\alph*.]
\item True or false---explain why: ``All relations are functions, but not all functions are relations.''
\item For each of the following, state whether \(y\) is a function of
  \(x\) (Y/N). If it is a function, state its domain. If it is not,
  state a restriction that would make it one.
  \begin{enumerate}[label=\roman*.]
  \item \(y = x^2\)
  \item \(x = y^2\)
  \item \(y = x^{1/2}\)
  \end{enumerate}
\item Is \(y = x^2\) invertible on all of \(\mathbb{R}\)? If not, give
  a domain on which it is invertible.
\end{enumerate}

\problem{Simplify expressions}

Simplify the following expressions as much as
possible:\footnote{inspired by Gill 1.1}

\begin{enumerate}[label=\alph*.]
\item \((-x^2y^3)^3\)
\item \((-2)^{(4-9)}\)
\item \(\left(\dfrac{1}{27b^4}\right)^{1/3}\)
\item \(\dfrac{13a/7b}{13b/2a}\)
\item \((a+b)^2 + (a-b)^2 + 2(a+b)(a-b) - 3a^2\)
\end{enumerate}

\problem{Graph sketching}

Let the functions \(f(x)\) and \(g(x)\) be defined for all
\(x \in \mathbb{R}\) by\footnote{inspired by Pemberton and Rau problem 3-1}
\[
f(x) = \left\{
    \begin{array}{ll}
        | x - 2 | & \quad \mbox{if} \ x < 1 \\
        1 & \quad \mbox{if} \ x \geq 1
    \end{array}
\right., \quad
g(x) = \left\{
    \begin{array}{ll}
        x^3 & \quad \mbox{if} \ x < 2 \\
        4 & \quad \mbox{if} \ x \geq 2
    \end{array}
\right.
\]

\begin{enumerate}[label=\alph*.]
\item Sketch \(y = f(x)\)
\item Compute \(f(g(0))\) and \(f(g(3))\)
\item Sketch \(y = f(g(x))\)
\item Sketch \(y = g(f(x))\)
\end{enumerate}

\problem{Root finding}

Find the roots (solutions) of the following quadratic
equations.\footnote{Gill 1.25} Recall that for \(ax^2 + bx + c = 0\),
\[
x = \dfrac{-b \pm \sqrt{b^2 - 4ac}}{2a}
\]

\begin{enumerate}[label=\alph*.]
\item \(9x^2 - 3x - 12 = 0\)
\item \(x^2 - 2x - 16 = 0\)
\item \(6x^2 - 6x - 6 = 0\)
\end{enumerate}

\problem{Simplify logarithms}

Express each of the following as a single logarithm:\footnote{Inspired
  by Grimmer HW1.3. Assume log is natural log unless specified
  otherwise.}

\begin{enumerate}[label=\alph*.]
\item \(\log(x) - 2\log(y) + \log(z)\)
\item \(2 \log(x) + \log(1)\)
\item \(\log(2x) - 2\)
\end{enumerate}

% Shared AI and Resources statement.
% Include in any problem set with:  % Shared AI and Resources statement.
% Include in any problem set with:  % Shared AI and Resources statement.
% Include in any problem set with:  \input{ai-statement}
\section*{AI and Resources statement}

Please list (in detail) all resources you used for this assignment. If
you worked with people, list them here as well. It is not enough to say
that you used a resource for help, you need to be specific on the link
and \emph{how} it was helpful. W/R/T gen AI tools (including GPT,
Claude, etc.) you cannot use them to do work on your behalf---you
cannot put in any of the questions, etc. You can ask for help on logic
/ sample problems. If you do use GPT or other AI tools, you need to
provide a link to your chat transcript. Any suspected academic
integrity violations will be immediately reported to the Dean of
Students.

\section*{AI and Resources statement}

Please list (in detail) all resources you used for this assignment. If
you worked with people, list them here as well. It is not enough to say
that you used a resource for help, you need to be specific on the link
and \emph{how} it was helpful. W/R/T gen AI tools (including GPT,
Claude, etc.) you cannot use them to do work on your behalf---you
cannot put in any of the questions, etc. You can ask for help on logic
/ sample problems. If you do use GPT or other AI tools, you need to
provide a link to your chat transcript. Any suspected academic
integrity violations will be immediately reported to the Dean of
Students.

\section*{AI and Resources statement}

Please list (in detail) all resources you used for this assignment. If
you worked with people, list them here as well. It is not enough to say
that you used a resource for help, you need to be specific on the link
and \emph{how} it was helpful. W/R/T gen AI tools (including GPT,
Claude, etc.) you cannot use them to do work on your behalf---you
cannot put in any of the questions, etc. You can ask for help on logic
/ sample problems. If you do use GPT or other AI tools, you need to
provide a link to your chat transcript. Any suspected academic
integrity violations will be immediately reported to the Dean of
Students.


\end{document}
